\section{KV Cache Compression with TurboQuant}
\label{sec:kv_cache}

\subsection{The KV Cache Bottleneck}

While weight quantization reduces model size, the KV cache grows with sequence length and becomes the dominant memory consumer for long contexts. For MiniMax M2.7 at full 192K context:

\begin{table}[h]
\centering
\caption{KV Cache Memory Requirements (192K context)}
\label{tab:kv_requirements}
\begin{tabular}{lcc}
\toprule
\textbf{Precision} & \textbf{Bytes/Token} & \textbf{Total (GB)} \\
\midrule
FP16 & 64 KB & 12.29 \\
Q8\_0 & 32 KB & 6.14 \\
Q4\_0 & 16 KB & 3.07 \\
Turbo4 (4-bit) & 18 KB & 3.46 \\
Turbo3 (3-bit) & 14 KB & 2.69 \\
\bottomrule
\end{tabular}
\end{table}

With Q4\_K\_M weights (1.35GB) and 2GB working memory, the combination of FP16 KV cache would require over 15GB---within our 24GB budget but leaving limited headroom. TurboQuant compression is essential for comfortable deployment.

\subsection{TurboQuant Methodology}

TurboQuant is llama.cpp's KV cache quantization mechanism, enabling aggressive compression while maintaining retrieval accuracy.

\subsubsection{Quantization Scheme}

TurboQuant applies per-channel quantization to keys and values:

\begin{equation}
Q(K) = \text{round}\left(\frac{K - z_K}{s_K}\right), \quad Q(V) = \text{round}\left(\frac{V - z_V}{s_V}\right)
\end{equation}

Where $s$ and $z$ are learned or computed per-head scales and zero-points.

\subsubsection{3-bit Implementation (Turbo3)}

The 3-bit implementation packs 10 values per 32-bit word with careful handling of outliers:

\begin{itemize}
    \item Standard 3-bit: 8 values per 32 bits
    \item Turbo3: 10 values per 32 bits (effective 3.2 bits)
    \item Special handling for extreme values to prevent overflow
\end{itemize}

\subsection{Accuracy Analysis}

We evaluate TurboQuant's impact on retrieval accuracy.

\subsubsection{Needle-in-Haystack Test}

The needle-in-haystack test places a specific fact at various positions in a long document and tests the model's ability to retrieve it.

\begin{table}[h]
\centering
\caption{TurboQuant Retrieval Accuracy}
\label{tab:turboquant_accuracy}
\begin{tabular}{lcccc}
\toprule
\textbf{Cache Type} & \textbf{4K} & \textbf{16K} & \textbf{64K} & \textbf{128K} \\
\midrule
FP16 & 100\% & 98\% & 96\% & 94\% \\
Q8\_0 & 100\% & 97\% & 95\% & 93\% \\
Turbo4 & 100\% & 96\% & 94\% & 91\% \\
Turbo3 & 98\% & 95\% & 92\% & 88\% \\
\bottomrule
\end{tabular}
\end{table}

Turbo3 maintains $>$90\% accuracy even at 128K context, sufficient for most applications.

\subsubsection{Depth-Dependent Analysis}

Retrieval accuracy varies by position within the context:

\begin{itemize}
    \item \textbf{Beginning (0-10\%):} Highest accuracy, 98-100\%
    \item \textbf{Middle (40-60\%):} Moderate accuracy, 92-95\%
    \item \textbf{End (90-100\%):} Good accuracy, 94-96\%
\end{itemize}

The ``lost in the middle'' phenomenon observed in some models is partially mitigated by MiniMax's architecture but still present.

\subsection{Perplexity Impact}

KV cache quantization affects perplexity differently than weight quantization:

\begin{table}[h]
\centering
\caption{Perplexity by KV Cache Quantization}
\label{tab:kv_perplexity}
\begin{tabular}{lcc}
\toprule
\textbf{KV Cache} & \textbf{Perplexity} & \textbf{Relative} \\
\midrule
FP16 & 12.45 & 1.00 \\
Q8\_0 & 12.51 & 1.005 \\
Turbo4 & 12.68 & 1.018 \\
Turbo3 & 12.89 & 1.035 \\
\bottomrule
\end{tabular}
\end{table}

The 3.5\% perplexity increase with Turbo3 is acceptable given the 5x memory reduction.

\subsection{Flash Attention Integration}

Flash Attention \cite{dao2022flashattention} reduces HBM accesses by fusing attention computation:

\begin{equation}
O = \text{softmax}\left(\frac{QK^T}{\sqrt{d_k}}\right)V
\end{equation}

Standard implementation:
\begin{enumerate}
    \item Load Q, K from HBM to SRAM
    \item Compute $S = QK^T$, store to HBM
    \item Load S, compute $P = \text{softmax}(S)$, store to HBM
    \item Load P, V, compute $O = PV$, store to HBM
\end{enumerate}

Flash Attention:
\begin{enumerate}
    \item Load Q, K, V tiles to SRAM
    \item Online softmax: compute in tiles, accumulate statistics
    \item Write only final O to HBM
\end{enumerate}

This reduces HBM accesses from $O(N^2)$ to $O(N)$, critical for memory-bandwidth-bound CPU inference.

\subsection{Context Length Strategy}

The MiniMax model supports 196,608 token context window. Our strategy:

\begin{equation}
n_{ctx} = \min(196608 - 2048, T_{max\_memory})
\end{equation}

We reserve 2,048 tokens for generation, yielding 194,560 usable context tokens (192K effective).

\subsubsection{Dynamic Context Sizing}

If memory is insufficient for full context, the system automatically reduces:

\begin{enumerate}
    \item Attempt 194,560
    \item If failure, try 131,072
    \item Continue halving: 65,536, 32,768, 16,384, 8,192, 4,096
    \item Minimum: 2,048 tokens
\end{enumerate}

This ensures the model loads even on severely constrained systems.

\subsection{Multi-Sequence Memory Sharing}

When processing multiple sequences concurrently, the KV cache can be shared for common prefixes:

\begin{itemize}
    \item System prompts shared across user requests
    \item Document contexts shared across questions
    \item Batch processing of similar inputs
\end{itemize}

This can reduce aggregate memory requirements by 20-40\% for multi-turn conversations.

\subsection{Cache Eviction Policies (Future Work)}

While Miniforge currently uses full-cache retention, we analyze potential eviction strategies:

\subsubsection{H2O (Heavy-Hitter Oracle)}

Retain tokens with highest accumulated attention scores:

\begin{equation}
\text{score}(t) = \sum_{l=1}^{L} \sum_{h=1}^{H} \text{attention}_{l,h}(t)
\end{equation}

This can reduce cache size by 50\% with minimal accuracy impact.

\subsubsection{StreamingLLM}

Maintain a fixed window of recent tokens plus ``attention sinks'' (initial tokens that receive disproportionate attention):

\begin{equation}
\text{Cache} = \text{Sinks} \cup \text{Recent}[t - w:t]
\end{equation}

This enables infinite context length with bounded memory, trading completeness for feasibility.

\subsection{Measurement and Monitoring}

Miniforge provides comprehensive KV cache monitoring through the memory management interface:

\textbf{Cache Registration:}
\begin{itemize}
    \item Cache allocations are tracked at initialization time
    \item Context length and compression type recorded
    \item Memory reserved and protected from overallocation
\end{itemize}

\textbf{Statistics Collection:}
\begin{itemize}
    \item KV cache allocation size (GB)
    \item Model weights allocation (GB)
    \item Available system memory (GB)
    \item Utilization percentage
\end{itemize}

These metrics enable real-time monitoring and debugging of memory consumption during inference operations.

\subsection{Performance Characteristics}

KV cache access patterns affect inference speed:

\begin{itemize}
    \item \textbf{Prompt processing:} Sequential cache writes, bandwidth-intensive
    \item \textbf{Token generation:} Random cache reads (attention over full context), latency-sensitive
    \item \textbf{Quantization overhead:} Dequantization adds 5-10\% compute overhead
\end{itemize}

TurboQuant's 3-bit format is optimized for fast dequantization on AVX2, minimizing overhead.
